% =====================================================================
%  Chapitre 1 - Etude prealable
% =====================================================================
\chapter{Etude prealable}
\label{chap:etude-prealable}

\section*{Introduction}
\addcontentsline{toc}{section}{Introduction}

Tout projet informatique repond a un besoin ne d'une situation concrete.
Ce premier chapitre etablit cette situation. Il presente l'organisme au
sein duquel le projet a ete conduit, formule la problematique a laquelle
il entend repondre, examine la maniere dont ce besoin est aujourd'hui
couvert et en releve les limites. Il expose ensuite la solution retenue
et justifie la methodologie de travail adoptee.

% =====================================================================
\section{Presentation de l'organisme d'accueil}

\subsection{L'entreprise}

Le present projet de fin d'etudes a ete realise au sein de
\essentiel{\entreprise}, entreprise de services numeriques implantee a
\entrepriseVille{} (\entrepriseCodePostal), en Tunisie.

L'activite de \entreprise{} s'organise autour de trois metiers
complementaires.

\begin{table}[H]
\centering
\caption{Les trois metiers de \entreprise}
\label{tab:metiers-sipro}
\small
\begin{tabularx}{\textwidth}{@{}p{3.6cm}X@{}}
\toprule
\textbf{Metier} & \textbf{Perimetre} \\ \midrule
Edition logicielle & Logiciels de gestion et progiciels d'entreprise.
  Metier historique, dont releve le present projet \\
Integration RFID & Identification par radiofrequence : tracabilite des
  biens, inventaire, controle d'acces \\
Materiel informatique & Vente, maintenance et assistance sur le materiel
  professionnel \\
\bottomrule
\end{tabularx}
\end{table}

Cette combinaison eclaire le projet. Une entreprise qui integre des
solutions de tracabilite et maintient du materiel sur site connait de
l'interieur les contraintes de la gestion d'un parc : suivi des
equipements, planification des interventions, surveillance des echeances.
\essentiel{Le sujet confie prolonge donc un savoir-faire existant vers un
domaine ou la demande des clients de \entreprise{} etait exprimee.}

\subsection{La mission confiee}

\entreprise{} a confie la conception et le developpement de
\essentiel{\sujetCourt}, une plateforme web complete de gestion de flotte
et de transport de marchandises, destinee a digitaliser les operations
des societes de transport clientes de l'entreprise.

Le stage s'est deroule du \stageDebut{} au \stageFin, soit une duree de
\stageDuree{} mois. Cette duree a rendu possible un decoupage en six
iterations d'un mois environ, chacune consacree a un domaine metier
coherent et livrant un increment utilisable.

Les objectifs fixes au depart etaient au nombre de quatre.

\begin{enumerate}
  \item \essentiel{Centraliser} dans un referentiel unique les vehicules,
    les equipements, les personnels et les clients, aujourd'hui disperses
    entre plusieurs documents.
  \item \essentiel{Industrialiser} la maintenance : ordres de travail,
    interventions, stock de pieces de rechange et surveillance des
    echeances reglementaires.
  \item \essentiel{Ouvrir un canal numerique au client final}, de la
    demande de transport au suivi de l'acheminement et a la consultation
    des factures.
  \item \essentiel{Outiller la decision} par l'estimation tarifaire,
    l'anticipation des pannes et l'affectation assistee des ressources.
\end{enumerate}

Le livrable attendu etait une plateforme fonctionnelle, deployable et
documentee, accompagnee de la modelisation UML correspondante.

% =====================================================================
\section{Problematique}
\label{sec:problematique}

Une entreprise de transport de marchandises pilote simultanement quatre
ressources dont les cycles de vie sont independants : un parc de
vehicules, un ensemble d'equipements embarques, une equipe de chauffeurs
et un portefeuille de clients. Chacune produit ses propres evenements,
et c'est de leur croisement que naissent les difficultes.

\figReelle{problematique.png}
  {Interdependance des ressources d'une entreprise de transport}
  {fig:problematique}{0.55}

La figure~\ref{fig:problematique} illustre cette interdependance. Elle se
traduit concretement par cinq difficultes recurrentes.

\paragraph{La dispersion de l'information.}
Le suivi des vehicules, celui des interventions de maintenance, le
planning des chauffeurs et la facturation vivent dans des documents
distincts, souvent detenus par des personnes differentes. Repondre a une
question aussi simple que \exemple{ce vehicule est-il disponible jeudi
prochain ?} suppose de consulter plusieurs sources et d'esperer qu'elles
soient a jour.

\paragraph{L'absence d'anticipation sur la maintenance.}
Les interventions sont declenchees par la panne plutot que par
l'observation de son approche. Le cout d'une immobilisation imprevue
depasse largement celui de l'entretien qui l'aurait evitee : il faut y
ajouter la mission non honoree, le report sur un autre vehicule et
l'insatisfaction du client.

\paragraph{Le risque administratif.}
Une visite technique expiree, une assurance echue ou un permis de
conduire perime exposent l'entreprise a une immobilisation immediate et a
des sanctions. Ces echeances sont nombreuses, dispersees dans le temps et
rarement surveillees de maniere systematique.

\paragraph{L'estimation tarifaire empirique.}
Le prix d'un transport se fixe souvent d'apres l'experience de la
personne qui etablit le devis. Deux devis comparables peuvent ainsi
differer sensiblement, ce qui fragilise la coherence commerciale et rend
la marge difficile a piloter.

\paragraph{L'opacite pour le client.}
Le client n'a pas de moyen autonome de deposer une demande, de suivre
l'acheminement de sa marchandise ni de consulter ses factures. Chaque
demande d'information mobilise un interlocuteur, ce qui alourdit le
travail administratif des deux cotes.

\vspace{0.4cm}
\noindent
\essentiel{Ces cinq difficultes ont une cause commune : l'absence d'un
referentiel unique ou les vehicules, les equipements, les chauffeurs, les
missions et les clients seraient decrits une seule fois et relies entre
eux.} C'est ce constat qui fonde le projet.

% =====================================================================
\section{Etude de l'existant}
\label{sec:existant}

\subsection{La situation avant le projet}

Chez les societes de transport clientes de \entreprise, la gestion
reposait sur des \essentiel{methodes traditionnelles}, articulees autour
de trois supports.

\begin{itemize}
  \item \essentiel{Des feuilles de calcul Excel} pour le suivi du parc,
    des missions et de la facturation. Chaque classeur etait tenu par une
    personne, sur son poste, dans sa propre organisation de colonnes.
  \item \essentiel{Des echanges par appel telephonique et par WhatsApp}
    pour la coordination quotidienne : affectation d'un chauffeur,
    signalement d'une panne, confirmation d'une livraison. Ces echanges
    portaient l'information operationnelle mais n'en laissaient aucune
    trace exploitable.
  \item \essentiel{Un tableau de bord rudimentaire}, alimente
    manuellement, offrant une vue partielle et rarement a jour.
\end{itemize}

Cette organisation fonctionne tant que le volume reste modeste et que les
memes personnes restent en poste. Elle atteint sa limite des que
plusieurs interlocuteurs doivent consulter ou modifier la meme
information, et surtout des qu'il faut \exemple{retrouver ce qui a ete
decide il y a trois mois}.

\subsection{Les alternatives du marche}

Face a cette situation, trois familles d'outils sont envisageables.

\paragraph{Le tableur.}
C'est la situation de depart decrite ci-dessus. Le tableur a pour lui sa
gratuite, sa souplesse et le fait que chacun sache s'en servir. Il ne
gere en revanche ni les acces concurrents, ni la tracabilite des
modifications, ni les regles metier.

\paragraph{Le progiciel de gestion generaliste.}
Certaines entreprises s'appuient sur un progiciel comptable ou commercial
auquel elles greffent le suivi de la flotte. La comptabilite et la
facturation y sont bien traitees, mais les notions propres au transport
(mission, echeance reglementaire, ordre de travail, piece de rechange)
n'existent pas nativement et doivent etre detournees d'autres objets.

\paragraph{La solution specialisee en ligne.}
Il existe des offres dediees a la gestion de flotte, generalement
proposees en abonnement. Elles couvrent bien le suivi des vehicules et la
geolocalisation, mais leur cout est proportionnel a la taille du parc,
leur parametrage suppose une adaptation aux usages du fournisseur, et
l'hebergement des donnees echappe a l'entreprise.

\subsection{Critique de l'existant}

Le tableau~\ref{tab:critique-existant} confronte ces trois familles aux
besoins identifies dans la section~\ref{sec:problematique}.

\begin{table}[H]
\centering
\caption{Couverture des besoins par les pratiques existantes}
\label{tab:critique-existant}
\small
\begin{tabularx}{\textwidth}{@{}Xccc@{}}
\toprule
\textbf{Besoin} & \textbf{Tableur} & \textbf{Progiciel} & \textbf{Solution} \\
                &                  & \textbf{generaliste} & \textbf{specialisee} \\
\midrule
Referentiel unique des ressources      & non      & partiel  & oui \\
Travail simultane a plusieurs          & non      & oui      & oui \\
Tracabilite des modifications          & non      & oui      & oui \\
Surveillance des echeances reglementaires & non   & non      & partiel \\
Maintenance preventive et ordres de travail & non & non      & partiel \\
Estimation tarifaire outillee          & non      & non      & non \\
Portail autonome pour le client        & non      & non      & partiel \\
Suivi de mission en direct             & non      & non      & oui \\
Maitrise des donnees et du cout        & oui      & partiel  & non \\
\bottomrule
\end{tabularx}
\end{table}

Aucune des trois familles ne couvre l'ensemble du besoin. Les deux
manques les plus nets concernent \essentiel{l'aide a la decision}
(estimation tarifaire, anticipation des pannes, affectation optimisee des
ressources), absente partout, et \essentiel{la surveillance automatique
des echeances}, qui n'est traitee qu'en partie et rarement de facon
transversale aux vehicules, aux chauffeurs et aux equipements.

\subsection{Solution proposee}

\sujetCourt{} est une plateforme unifiee de gestion de flotte et de
transport de marchandises. Elle repose sur trois principes.

\paragraph{Un referentiel unique, des domaines separes.}
Chaque ressource est decrite une seule fois. Les domaines metier restent
toutefois autonomes : la facturation, la maintenance et l'exploitation
evoluent independamment, ce qui permet de faire progresser un domaine
sans immobiliser les autres.

\paragraph{Trois publics, trois applications.}
Un back-office pour le personnel de l'entreprise, un portail pour les
clients, une application mobile pour les chauffeurs. Chacune n'expose que
ce qui concerne son public.

\paragraph{L'aide a la decision integree au flux de travail.}
L'estimation tarifaire, l'anticipation des pannes, l'affectation des
ressources et la veille metier ne sont pas des modules a part : elles
interviennent au moment ou la decision se prend, dans l'ecran ou
l'utilisateur travaille.

\begin{definitionbox}{Ce que couvre la plateforme}
Douze services metier autonomes, trois applications destinees a des
publics distincts, un referentiel commun de vehicules, d'equipements,
de personnels et de clients, et quatre services d'aide a la decision.
Le detail de ce perimetre est etabli au chapitre suivant.
\end{definitionbox}

% =====================================================================
\section{Methodologie de travail}
\label{sec:methodologie}

Le choix d'une methodologie n'est pas neutre : il determine la maniere
dont les besoins sont recueillis, dont le travail est decoupe et dont
l'avancement est mesure.

\subsection{Le developpement en cascade}

Le modele en cascade enchaine des phases sequentielles : specification,
conception, developpement, tests, mise en service. Chaque phase se
termine avant que la suivante ne commence, et produit un document qui
sert de contrat pour la suite.

Ce modele convient aux projets dont le besoin est stable et connu des le
depart. Il presente en revanche deux inconvenients pour un projet comme
celui-ci : le client ne voit un resultat utilisable qu'a la toute fin, et
une erreur de comprehension du besoin ne se decouvre qu'au moment des
tests, lorsque son cout de correction est maximal.

\subsection{L'approche agile}

Les approches agiles renversent cette logique. Le travail est decoupe en
iterations courtes, chacune produisant un increment fonctionnel du
produit. Le besoin est precise au fil de l'eau, au contact du resultat
plutot que sur document.

\subsection{Le cadre Scrum}

Scrum est le cadre agile le plus repandu. Il repose sur trois roles, sur
des evenements qui rythment l'iteration et sur trois artefacts.

\figReelle{methodo-scrum.jpg}
  {Le processus Scrum}
  {fig:scrum}{0.72}

\begin{table}[H]
\centering
\caption{Roles, evenements et artefacts de Scrum}
\label{tab:scrum}
\small
\begin{tabularx}{\textwidth}{@{}p{2.8cm}X@{}}
\toprule
\textbf{Element} & \textbf{Definition} \\ \midrule
\textit{Roles} & Product Owner : porte la valeur metier et ordonne le
  carnet de produit. Scrum Master : garantit le cadre et leve les
  obstacles. Equipe de developpement : concoit, developpe et teste
  l'increment \\
\textit{Evenements} & Planification, melee quotidienne, revue de sprint,
  retrospective \\
\textit{Artefacts} & Carnet de produit, carnet de sprint, increment \\
\bottomrule
\end{tabularx}
\end{table}

\subsection{Choix retenu et adaptation au contexte}

\essentiel{Scrum a ete retenu pour ce projet}, pour trois raisons.

D'abord, \essentiel{le besoin n'etait pas entierement connu au depart}.
Le perimetre s'est precise au contact des utilisateurs : plusieurs
fonctionnalites, comme l'archivage reversible des comptes ou la
surveillance des echeances documentaires, ont ete identifiees en cours de
projet et n'auraient pas figure dans une specification initiale.

Ensuite, \essentiel{le decoupage en domaines metier se prete
naturellement a l'iteration}. Chaque sprint peut livrer un domaine
complet et utilisable, ce qui rend l'avancement lisible et permet un
retour utilisateur rapide.

Enfin, \essentiel{le format du projet impose un rythme}. Un stage de fin
d'etudes se deroule sur une duree fixe : le decoupage en iterations de
duree egale donne une mesure objective de l'avancement et permet de
reajuster le perimetre en cours de route plutot qu'a la fin.

\paragraph{Adaptation au contexte du stage.}
Scrum suppose une equipe de plusieurs developpeurs. Le projet ayant ete
mene par un seul developpeur, le cadre a ete adapte : les roles de
Product Owner et de Scrum Master ont ete tenus par les encadrants, la
melee quotidienne est devenue un point d'avancement hebdomadaire, et la
revue de sprint a pris la forme d'une demonstration a l'encadrant
entreprise. Les artefacts, en revanche, ont ete tenus integralement :
carnet de produit, carnet de sprint et increment livre a chaque
iteration.

\subsection{Outils de pilotage}
\label{sec:outils-pilotage}

Trois outils ont soutenu la conduite du projet au quotidien.

\begin{table}[H]
\centering
\caption{Outils de pilotage du projet}
\label{tab:outils-pilotage}
\small
\begin{tabularx}{\textwidth}{@{}p{2.2cm}p{3.4cm}X@{}}
\toprule
\textbf{Outil} & \textbf{Usage} & \textbf{Apport} \\ \midrule
Trello & Carnet et suivi des taches & Chaque recit devient une carte,
  deplacee au fil de son avancement : le carnet de sprint est visible
  d'un coup d'oeil \\
Git et GitHub & Gestion de versions & Historique, branches, retour
  arriere, sauvegarde et support aux revues de code \\
WhatsApp & Communication quotidienne & Canal court avec l'encadrant
  entreprise, sans attendre le point hebdomadaire \\
\bottomrule
\end{tabularx}
\end{table}

\begin{figure}[H]
  \centering
  \includegraphics[height=1.35cm,keepaspectratio]{trello-logo.jpg}\hspace{1.4cm}
  \includegraphics[height=1.35cm,keepaspectratio]{logo-git.png}\hspace{1.4cm}
  \includegraphics[height=1.35cm,keepaspectratio]{whatsapplogo.jpg}
  \caption{Outils de pilotage et de communication du projet}
  \label{fig:outils-pilotage}
\end{figure}

\essentiel{WhatsApp figure a la fois dans l'existant critique et dans les
outils de pilotage.} La contradiction n'est qu'apparente : ce qui pose
probleme n'est pas l'outil mais son emploi comme \exemple{support
d'information metier}. Coordonner une equipe par messagerie instantanee
est legitime ; y conserver l'etat d'un parc de vehicules ne l'est pas.
C'est la frontiere que la plateforme etablit.

% =====================================================================
\section*{Conclusion}
\addcontentsline{toc}{section}{Conclusion}

Ce chapitre a etabli le contexte du projet. L'examen des pratiques
existantes a montre qu'aucune des trois familles d'outils rencontrees ne
couvre l'ensemble du besoin d'une entreprise de transport de
marchandises, les manques les plus nets portant sur l'aide a la decision
et sur la surveillance des echeances reglementaires. La solution
proposee, \sujetCourt, repose sur un referentiel unique servant des
domaines metier autonomes, sur trois applications adaptees a trois
publics et sur une aide a la decision integree au flux de travail. Le
cadre Scrum a ete retenu pour conduire cette realisation, adapte au
contexte d'un stage, et outille par Trello, Git et GitHub.

Le chapitre suivant traduit ce cadrage en exigences : il identifie les
acteurs du systeme, recense les besoins fonctionnels et non fonctionnels,
et etablit le decoupage du travail en six sprints.
